\section{詳細解説}

\subsection{円錐の軸の向きと回転の大きさ}

人間大の 3 次元モデルで、半頂角 30° の円錐の軸の向き $(\psi_0, \alpha_0)$ を変えて、1 周あたりの胴体の回転の大きさを調べた（図~\ref{fig:D10_cone_axes.png}）。
\begin{itemize}
\item 回転の大きさは 5.1〜38.2° の範囲で、軸の向きによって 7 倍以上変わる。
\item 軸が水平（$\alpha_0=0$）に近いと小さく、上下（$\alpha_0=\pm60^\circ$）に近いと大きい。
  軸が鉛直に近いと胴体は縦軸まわり（ヨー）に回り、胴体の縦軸まわりの慣性が最も小さい（細長い体）ためである。
\item 最大は軸が左下方（$\psi_0=90^\circ$, $\alpha_0=-60^\circ$）のときで、右肩から見て体の内側・下向きに腕を回す動きにあたる。
\end{itemize}
\textbf{どの軸まわりに向きを変えたいかによって、効率の良い腕の回し方が大きく異なる}。3 次元の AMBAC の計画では、この地図が「どの動作を選ぶか」の基礎になる。

\DFIG{D10_cone_axes.png}{円錐の軸の向き（横軸：方位角、縦軸：仰角）と、半頂角 30° の円錐 1 周あたりの胴体の回転 [°]。}

\subsection{非可換性の数値例}

ヨー用の円錐（軸：真上）1 周で約 58°、ピッチ用の円錐（軸：右横）3 周で約 33° 回るとき、
\begin{itemize}
\item ヨー → ピッチの最終姿勢：回転ベクトル $(8.2^\circ, 41.4^\circ, -64.1^\circ)$
\item ピッチ → ヨーの最終姿勢：回転ベクトル $(-22.4^\circ, 33.7^\circ, -65.3^\circ)$
\end{itemize}
で、2 つの姿勢の差は 29.4° だった。回転が小さければ順序の影響は 2 次の小ささ（角度の積に比例）だが、数十度の回転では無視できない。

\subsection{ジンバルロックの数値例}

腕先を一定の角速度 1.57 rad/s で動かす円錐が、極（真上）から角距離 $\delta$ のところを通るとき、方位角の関節速度の最大値は
付録~\ref{app:3d} の見積もり $\dot\psi\approx\omega_{\text{tip}}/\cos\alpha\approx\omega_{\text{tip}}/\delta$ どおりになった。
\begin{longtable}{rrr}
\toprule
$\delta$ & 見積もり $1.57/\delta$ & シミュレーション \\ \midrule \endhead
10° & 9.0 rad/s & 9.05 rad/s \\
3° & 30.0 rad/s & 30.0 rad/s \\
1° & 90.0 rad/s & 90.0 rad/s \\
\bottomrule
\end{longtable}

\subsection{定格荷重マップの数値}

\DTAB{sim07_load_chart}

方位 90°（真横）で最も小さく、仰角を上げるほど大きくなる。仰角 60° 以上では、荷物が接地面の真上に来るため、静的にはいくら重くても転倒しない。

\subsection{旋回の最適化の定式化}

180° の水平旋回では、方位角 $\psi(\tau)$ を最小躍度のプロファイル＋3 つの補正、仰角 $\alpha(\tau)$ を $\sin^2(\pi\tau)$ を基本とする 4 つの基底関数で表し、
動作時間 $T$ を最小化した。制約は、転倒（接地長方形の内側 5 cm 以内）、滑り、関節の速度 6 rad/s・加速度 40 rad/s$^2$、仰角 $-35^\circ$ 以上 85° 以下である。
初期値に「水平のまま」「60° 持ち上げる」「35° 下げる」の 3 通りとそれらにランダムな補正を加えたものを使い、局所解を避けた。

\section*{考察課題}
\begin{exercise}
3 次元で目標の姿勢に最短で到達するには、図~\ref{fig:D10_cone_axes.png} のような「動作ごとの回転の地図」をどう使えばよいか。非可換性を考慮した計画の手順を考えよ。
\end{exercise}
\begin{exercise}
肩の 2 軸ジンバルの取り付け方向を変えて、極を真上ではなく真後ろに置くと、どの動作で特異点の問題が起きるか。
\end{exercise}
